\documentclass{cumcmthesis}
\usepackage[framemethod=TikZ]{mdframed}
\usepackage{url}   % 网页链接
\usepackage{subcaption} % 子标题
\usepackage{booktabs} % 浮动体
\usepackage{amsmath}


\title{我的论文标题}
\tihao{C}
\baominghao{121}
\schoolname{上海交通大学}
\membera{张航}
\memberb{刘允禾}
\memberc{高慧鑫}
\supervisor{高晓沨}
\yearinput{2026}     % 年
\monthinput{09}      % 月
\dayinput{13}

\begin{document}

\maketitle

\begin{abstract}
            
            \keywords{关键词1\quad  关键词2\quad   关键词3}
        \end{abstract}
        \section{问题重述}
        \subsection{问题背景}
        \subsection{问题重述}
        \section{问题分析}
        \subsection{问题一的分析}
        \subsection{问题二的分析}
        \subsection{问题三的分析}
        \section{模型假设}
        \section{符号说明}
        \section{数据预处理}
        \begin{center}
            \begin{tabular}{cc}
                \hline
                \makebox[0.3\textwidth][c]{符号}	&  \makebox[0.4\textwidth][c]{意义} \\ \hline
                D	    & 木条宽度（cm） \\ \hline
            \end{tabular}
        \end{center}
        \section{问题一模型的建立与求解}
        \subsection{模型建立}
        \subsection{约束条件}
        \subsection{模型求解}
        \section{问题二模型的建立与求解} 
        \subsection{模型建立}
        \subsection{约束条件} 
        \subsection{模型求解}
        \section{问题三模型的建立与求解} 
        \subsection{模型建立}
        \subsection{约束条件} 
        \subsection{模型求解}
        \begin{equation}\tag{8}
            \sum_{j=16}^{\infty} x_{ijt} + \frac{\sum_{n=1}^{\infty} x_{ijt}}{2} \leq 1,\quad 27\leq i\leq 34
        \end{equation}
        \section{模型的评价、改进与推广}

        \begin{thebibliography}{9}%宽度9
            \bibitem{bib:one} ....
        \end{thebibliography}

        \begin{appendices}
            附录的内容。
        \end{appendices}
    \end{document}
\end{document}